\documentclass[12pt]{article}

% Import preambles and macros for homework
\input{../../../../preambles/hw-preamble.tex}
\input{../../../../preambles/homework-macros.tex}
\input{../../../../preambles/homework-title.tex}
\input{../../../../preambles/homework-config.tex}

\renewcommand{\author}{Deepak Jassal}
\renewcommand{\authorlast}{Jassal}
\renewcommand{\coursename}{Advanced Linear Algebra}
\renewcommand{\coursecode}{MATH 326}
\renewcommand{\assignment}{Assignment 7}
\renewcommand{\instructor}{Dr. Edward Dobrowolski}
\renewcommand{\duedate}{April 13\textsuperscript{th}, 2026}
\usepackage{tensor}

\begin{document}
\begin{titlepage}
	\centering
	\vspace*{2.0cm}	
	\pgfornament{84}\\
	{\LARGE \textsc{\coursename}\par}
	\vspace{0.5cm}
	{\large\coursecode\par}
    \vspace{0.5cm}
    {\large\instructor\par}
	\vspace{1.5cm}
	{\huge\bfseries\assignment\par}
	\vspace{1cm}  
	{\LARGE\itshape\author\par}
    \vspace{2cm}
	{\large\bfseries Due Date:\par}
	\vspace{0.5cm}
	{\Large \duedate}\\
	\pgfornament{84}
\end{titlepage}
\stepcounter{section}
\section*{Problem 1}
Let
\[
    A=
    \begin{bmatrix}
        2&1&1&1\\
        1&3&3&1\\
        1&0&1&2\\
        0&1&1&3
    \end{bmatrix}.
\]
Find a (psuedo) rational form of $A$. That is find a matrix $B$ such that $B^{-1}AB$ has the form
\[
    B^{-1}AB=
    \begin{bmatrix}
        0&0&0&-c_0\\
        1&0&0&-c_1\\
        0&1&0&-c_2\\
        0&0&1&-c_3
    \end{bmatrix},
\]
and explain what $c_0,c_1,c_2,c_3$ are. Software is allowed.\\
\textit{Solution.} The $c_i$ are the coefficients of the characteristic polynomial of $A$. Specifically, if $p(\lambda) = \det(\lambda I-A)$ is written as
\[
    p(\lambda)=\lambda^4-c_3\lambda^3+c_2\lambda^2-c_1\lambda+c_0,
\]
the last column of $BAB^{-1}$ is $(-c_0,-c_1,-c_2,-c_3)^T$. This is the companion matrtix of $p(\lambda)$, and it works because by the Cayley-Hamilton theorem $p(A)=0$, so the relation $A^4b=c_0b+c_1Ab+c_2A^2b+c_3A^3b$ holds for any vector $b$.\\
Compute the characteristic polynomial we have. 
\[
    p(\lambda)=\lambda^4-9\lambda^3+24\lambda^2-27\lambda+14.
\]
Comparing the signs we have
\[
    c_3=-9,c_2=24,c_1=-27,c_0=14.
\]
\[
    \begin{bmatrix}
        0&0&0&-14\\
        1&0&0&27\\
        0&1&0&-24\\
        0&0&1&27
    \end{bmatrix}.
\]
The matrix
\[
    B=
    \begin{bmatrix}
        1&2&6&25\\
        1&1&8&41\\
        0&1&3&13\\
        0&0&2&17
    \end{bmatrix}
\]
is a solution to $BAB^{-1}$.
\stepcounter{section}
\section*{Problem 2}
Let
\[
    A=
    \begin{bmatrix}
        11&-1&-5&-1\\
        -1&11&-1&-5\\
        -5&-1&11&-1\\
        -1&-5&-1&11
    \end{bmatrix}.
\]
Find an orthogonal matrix $Q$ such that $Q^TAQ$ is diagonal. Use of sotware is recommended.\\
\textit{Solution.} I found that
\[
    Q=
    \begin{bmatrix}
        \frac{1}{2}&-\frac{1}{2}&-\frac{1}{\sqrt{2}}&0\\
        \frac{1}{2}&\frac{1}{2}&0&-\frac{1}{\sqrt{2}}\\
        \frac{1}{2}&-\frac{1}{2}&-\frac{1}{\sqrt{2}}&0\\
        \frac{1}{2}&\frac{1}{2}&0&-\frac{1}{\sqrt{2}}
    \end{bmatrix}.
\]
Works, with $Q^TQ=I$ and
\[
    Q^TAQ=
    \begin{bmatrix}
        4&0&0&0\\
        0&8&0&0\\
        0&0&16&0\\
        0&0&0&16
    \end{bmatrix}.
\]

\stepcounter{section}
\section*{Problem 3}
Determine if
\[
    A=
    \begin{bmatrix}
        1&1&0\\
        0&1&1\\
        1&0&1
    \end{bmatrix}
\]
is unitarily diagonalizable.\\
\textit{Solution.} Note that 
\[
    AA^\ast=
    \begin{bmatrix}
        2&1&1\\
        1&2&1\\
        1&1&2
    \end{bmatrix}
\]
and
\[
    A^\ast A=
    \begin{bmatrix}
        2&1&1\\
        1&2&1\\
        1&1&2
    \end{bmatrix}.
\]
Since $AA^\ast=A^\ast A$, we have that $A$ is unitarily diagonalizable. 

\stepcounter{section}
\section*{Problem 4}
Prove that real skew-symmetric matrices are unitarily diagonalizable. Recall $A\in M_n(\R)$ is called skew-symmetric if $A^T=-A$.
\begin{proof}
    Let $A\in M_n(\R)$ be skew symmetric. Note that since $A$ is real we have $A^\ast=A^T=-A$. Thus we have
    \begin{align*}
        AA^\ast&=A^\ast A\\
        AA^T&=A^TA\\
        -A^2&=-A^2.
    \end{align*}
    Thus, we have that all real skew symmetric matrices are unitarily diagonalizable.
\end{proof}

\end{document}